\documentclass[CMPE]{KGCOEReport}

\usepackage{graphicx}
\usepackage{newtxtext, newtxmath}
\usepackage{csquotes}
\usepackage{parskip}
\usepackage{karnaugh-map}
\usepackage{float}
\usepackage{hyperref}
\usepackage[super]{nth}
\usepackage{pdfpages}

\graphicspath{{./assets/}}

\newcommand{\classCode}{CMPE 260}
\newcommand{\name}{Aden Perry}
\newcommand{\LabSectionNum}{L2}
\newcommand{\LabInstructor}{Richard Cliver}

\newcommand{\TAs}{Eric Reed \\ Tevin Hendess}
\newcommand{\LectureSectionNum}{01}
\newcommand{\LectureInstructor}{Cory Merkel}
\newcommand{\exerciseNumber}{7}
\newcommand{\exerciseDescription}{Processor Timing}
\newcommand{\dateDone}{2026-04-21}
\newcommand{\dateSubmitted}{2026-04-28}

\begin{document}

\maketitle

\section*{Questions}
\addcontentsline{toc}{section}{Questions}

What are setup and hold times?

Setup time is the time before the clock edge that the data signal must stay
stable for the flip-flop output to be trusted. Hold time is the time after the
clock edge that the data signal must stay stable for the flip-flop edge to be
trusted.

How would violating setup and hold times cause a design to fail?

Violating either can cause unpredictable or signal outputs after a clock edge.

\section*{Results and Analysis}
\addcontentsline{toc}{section}{Results and Analysis}

\begin{table}[H]
    \center
	\caption{Results}
	\begin{tabular}{ccc}
		Clock Frequency (MHz) & Worst Negative Slack (ns) & Worst Hold Slack (ns) \\
		\hline
		\hline
		10                    & 42.729                    & 0.028                 \\
		15                    & 24.481                    & 0.009                 \\
		20                    & 17.249                    & 0.057                 \\
		25                    & 12.158                    & 0.062                 \\
		30                    & 8.921                     & 0.016                 \\
		35                    & 6.907                     & 0.020                 \\
		40                    & 5.233                     & 0.035                 \\
		45                    & 2.960                     & 0.012                 \\
		50                    & 1.252                     & 0.038                 \\
		55                    & 1.041                     & 0.032                 \\
		60                    & 0.749                     & 0.026                 \\
		65                    & 0.535                     & 0.023                 \\
		70                    & 0.279                     & 0.036                 \\
		75                    & 0.547                     & 0.041                 \\
		80                    & 0.307                     & 0.015                 \\
		81                    & -0.658                    & 0.045                 \\
		85                    & -0.227                    & 0.025                 \\
		\hline
		\hline
	\end{tabular}
	\label{tab}
\end{table}

Interestingly, at high clock frequencies, WNS does not always go down with
increased frequency. Based on Table \ref{tab}, the highest clock frequency this
MIPS processor can support is 80 MHz.

\section*{Conclusion}
\addcontentsline{toc}{section}{Conclusion}

The requested output frequency of the clock of the MIPS design was incremented
by 5 MHz repeatedly, until a negative WNS was attained. Then, starting at the
last frequency to give a positive WNS, the frequency was incremented by 1 ns.
Because changing and re-implementing the design in Vivado is a time-consuming
process, the first 80 MHz tests were automated using a basic Tcl script. The
maximum working clock frequency was determined; the exercise was successful.
Tcl is a common language encountered in chip design software; basic knowledge
of Tcl will be useful in future projects.

\includepdf[pages=-]{assets/prelab}
\end{document}
